\section{References}

Each entry names the edition actually used and the exact loci this study read.
Where an artifact was acquired and hashed, its role and rights status are given
in the repository's source library rather than repeated here.

\subsection*{Controlling liturgical text}

\textbf{Missale Romanum ex decreto Sacrosancti Concilii Tridentini
restitutum, Summorum Pontificum cura recognitum}, editio typica (Vatican City:
Typis Polyglottis Vaticanis, 1962). Read in the 1{,}088-page facsimile hosted
by the Church Music Association of America, acquired and hashed before use;
exact bytes recorded, not retained. Loci read: \work{Rubricae generales
Missalis romani} nn.~424--432 and 466--514 (printed pp.~XXX--XXXV, artifact
pp.~28--33); \work{Ritus servandus in celebratione Missae} V--VIII (printed
pp.~LVII--LX, artifact pp.~59--62); \work{De defectibus} V (printed p.~LXVII,
artifact p.~69); \work{Ordo Missae} nn.~1013--1042 (printed pp.~216--222,
artifact pp.~297--303); \work{Praefationes} nn.~1043--1087 (printed
pp.~223--298, artifact pp.~304--379); \work{Canon Missae} nn.~1088--1112
(printed pp.~299--311, artifact pp.~380--392); \work{Ordo Missae}
nn.~1113--1139 (printed pp.~312--328, artifact pp.~393--409); and the front
matter carrying \latincue{Quo primum} (14 July 1570), the breve of Clement VIII
(7 July 1604), the breve of Urban VIII (2 September 1634), the general decree
of the Sacred Congregation of Rites on the code of rubrics (26 July 1960) and
the decree declaring this edition typical (23 June 1962).

\subsection*{English witness for the received prayers}

\textbf{The Roman Missal translated into the English language for the use of
the laity}, first revised edition (Philadelphia: Eugene Cummiskey, 1861). Read
in the Internet Archive scan of item \texttt{romanmissaltran00churgoog},
acquired and hashed before use; exact bytes recorded, not retained. Loci read:
printed pp.~xv--xix (the preparation, \latincue{Kyrie} and \latincue{Gloria});
pp.~xx--xxvi (lessons, Creed, offertory); p.~xxviii (the common preface and
\latincue{Sanctus}); pp.~xxxiii--xlii (the Canon, the Lord's Prayer, the
fraction, the \latincue{Agnus Dei} and the priest's Communion);
pp.~xliii--xlvi (the ablutions, the dismissal, \latincue{Placeat}, the blessing
and the last Gospel). A pre-1955 lay hand missal, not an official or approved
liturgical translation; quoted as a historical witness only.

\subsection*{Scripture}

\textbf{The Holy Bible, Douay-Rheims version}, Challoner revision, in the
electronic text distributed as Project Gutenberg eBook 1581. Verses quoted:
Genesis 4, 4 and 14, 18; Psalms 17, 4; 25, 6; 42, 4; 115, 12--13; 117, 26;
Isaias 6, 3 and 6--7; Daniel 3, 39--40; Matthew 6, 11; 8, 8; 21, 9; 26,
26--28; Luke 1, 11; 2, 14; 11, 3; John 1, 14 and 1, 29; 14, 27; Romans 12, 1;
1 Corinthians 11, 25 and 11, 29; Apocalypse 8, 3. Psalm numbering is the
Vulgate numbering the Missal prints.

\subsection*{Conciliar and papal acts}

\textbf{Second Vatican Council, \work{Lumen gentium} 10--11}, official English
dated web edition registered on 23 July 2026. The complete paragraphs control
the ordered relation and essential distinction between common and ministerial
priesthood and the faithful's Eucharistic participation.

\textbf{Canones et decreta sacrosancti oecumenici Concilii Tridentini}, editio
stereotypa undecima (Lipsiae: ex officina Bernhardi Tauchnitz, 1887, reprinting
the Roman edition of 1834). Read at the page images of printed pp.~118--121:
session XXII chapters III, IV, V and the end of VII, and canons I--II
and III--IX.

\textbf{The Canons and Decrees of the Sacred and Oecumenical Council of
Trent}, trans.\ J.~Waterworth (London 1848; Burns and Oates reprint). Read at
the page images of printed pp.~156 and 159: chapters IV and V, and canons
III--IX.

\textbf{Acta Apostolicae Sedis} 54 (1962), p.~873: Sacra Congregatio Rituum,
decree \latincue{De S. Ioseph nomine Canoni Missae inserendo}, 13 November
1962, in force from 8 December 1962.

\textbf{Acta Apostolicae Sedis} 52 (1960), pp.~593--596: John XXIII, motu
proprio \latincue{Rubricarum instructum}, and the general decree of the Sacred
Congregation of Rites promulgating the new code of rubrics. Read in the
volume's optical text layer; not collated at the page images, and cited here
only for the acts' identity and dates, which the 1962 Missal's own front matter
independently supplies.

\subsection*{Patristic witness}

\textbf{Patrologia Latina} tomus XVI, \emph{S. Ambrosius} (Parisiis: apud
Garnier Fratres et J.-P. Migne successores, 1880), coll.~462--464:
\work{De sacramentis} IV, 5, 21--23 and IV, 6, 27, with the Maurist notes on
the same columns. Read at the page images. The attribution of \work{De
sacramentis} to Ambrose is disputed and this study takes no position on it;
the text is used as a dated Milanese witness of about 390, which is all the
argument requires.

\subsection*{Modern scholarship}

Adrian Fortescue, \textbf{The Mass: A Study of the Roman Liturgy} (London and
New York: Longmans, Green and Co., 1922). Read at the page images of printed
pp.~172 (Gregory the Great and the Canon; Benedict XIV's dictum), 227 (the
prayers at the foot of the altar; the medieval uses; the papal chapel and
\latincue{sancta sanctorum}), 232 and 234 (the \latincue{Kyrie} and Gregory's
letter to John of Syracuse), 288 (the Creed at Rome in 1014, with the
dissenting authorities), 305 (the offertory prayers, the mixed chalice, and
\latincue{immaculatam hostiam} as ``an anticipation of the consecration''), 387
(the \latincue{Agnus Dei} and Sergius I) and 394 (the last Gospel before and
after 1570). Fortescue died in 1923 and this printing carries no revision by a
later hand; his own footnotes preserve the disagreements among the authorities
he cites, and those disagreements are reported here rather than resolved. His
datings rest on manuscript arguments this study has not examined and are
attributed to him throughout.
